\documentclass[11pt,a4paper]{article}
\usepackage[T1]{fontenc}
\usepackage[utf8]{inputenc}
\usepackage[margin=2.3cm]{geometry}
\usepackage{lmodern}
\usepackage{amsmath}
\usepackage{booktabs}
\usepackage{array}
\usepackage{eurosym}
\usepackage[hidelinks]{hyperref}
\usepackage{microtype}
\setlength{\parindent}{0pt}
\setlength{\parskip}{0.55em}

\title{Estimating the growth of the proton-therapy market\\
\large Treatment-room scenarios for 2036 and 2046}
\author{Working note for the ARGOS investment discussion}
\date{13 September 2026}

\begin{document}
\maketitle

\section{Purpose and preliminary conclusion}

This note develops an initial estimate of the worldwide proton-therapy
installed base over the next ten and twenty years. It is intended to support
discussion of the potential market for ARGOS, rather than to forecast ARGOS
sales directly.

The proposed central case is approximately \textbf{500 operating proton-treatment
rooms in 2036 and 800 in 2046}, starting from a \emph{provisional} baseline of
300 rooms in 2026 and assuming 5\% annual net growth. Slower and faster growth
cases use 3\% and 7\%, respectively. The baseline room count requires verification
before these figures are presented as an evidence-based market estimate.

These are scenario calculations constructed for this discussion, not published
market forecasts or statistical confidence intervals.

\section{What is being counted?}

The unit of analysis is an \textbf{operating clinical proton-treatment room}.
A centre can contain several rooms, and an accelerator can supply more than
one room. Consequently, centres, accelerators, room orders and operating rooms
must not be treated as interchangeable quantities.

The present ARGOS commercial model assumes one monitoring system per treatment
room. This is a commercial and technical assumption, not an established mapping
for every facility. Actual eligibility depends on room geometry, patient
transfer arrangements, gantry or upright treatment configuration and clinical
workflow. Carbon-ion-only rooms and research-only rooms should not be included
in a proton-room inventory. Dedicated ocular rooms should be identified
separately when assessing ARGOS compatibility.

\section{Evidence available for the starting point and pipeline}

\begin{itemize}
\item A June 2026 publication compiling PTCOG data reports \textbf{128
proton-therapy facilities operating in January 2026}. It also reports that
approximately 65\% of facilities under construction or planned have a single
treatment room. The increasing share of single-room projects means that centre
growth need not equal room growth~\cite{yap2026}.

\item A manual tally of PTCOG's June 2026 construction list gives approximately
\textbf{57 proton-treatment rooms across 32 proton-capable projects}, counting
only the proton rooms in the mixed proton/carbon project and excluding
carbon-only projects. Some listed completion dates are old or uncertain.
This is a tally of the published list, not independent confirmation that all
projects remain active or will open on schedule~\cite{ptcogconstruction}.

\item IBA reported \textbf{43 equipment projects in its year-end 2025 backlog}:
36 single-room ProteusONE projects and seven multi-room ProteusPLUS projects.
This manufacturer backlog overlaps the facility pipeline and must not be added
to the PTCOG count as a separate set of installations~\cite{iba2025}.
\end{itemize}

The evidence establishes an existing clinical network and a substantial
development pipeline. It does not by itself determine a worldwide treatment-room
count or a twenty-year growth rate. The \textbf{300-room baseline remains a
planning assumption}; it has not been obtained here from a reconciled,
facility-by-facility inventory. In particular, it should not be described as
300 verified ARGOS-compatible rooms.

\section{Growth model}

Let $N_0$ be the number of operating proton-treatment rooms at the 2026 reference
point and $g$ the annual net growth rate. The scenario model is
\begin{equation}
N(t)=N_0(1+g)^t,\qquad N_0=300,
\end{equation}
where $t$ is the number of years after 2026. Thus $t=10$ corresponds to 2036 and
$t=20$ to 2046. Results are rounded to whole rooms for display; calculations
should retain full precision internally.

The three annual growth assumptions represent different development paths:
\begin{description}
\item[Slow growth: 3\%.] Continued expansion constrained by capital requirements,
reimbursement, staffing and project execution.
\item[Central case: 5\%.] Steady deployment of compact systems, geographic
expansion and gradual improvement in clinical acceptance.
\item[Fast growth: 7\%.] Stronger adoption supported by favourable clinical
evidence, reimbursement and lower deployment costs.
\end{description}

These descriptions are modelling judgements. The rates have not been fitted to
a historical room-count series. They also assume a constant rate over each
horizon; actual expansion is likely to be uneven.

\begin{table}[htbp]
\centering
\begin{tabular}{lrrrr}
\toprule
Scenario & Annual growth & 2026 rooms & 2036 rooms & 2046 rooms\\
\midrule
Slow & 3\% & 300 & 403 & 542\\
Central & 5\% & 300 & 489 & 796\\
Fast & 7\% & 300 & 590 & 1,161\\
\bottomrule
\end{tabular}
\caption{Worldwide operating proton-treatment rooms under three assumed growth
rates. Figures are calculated scenarios, not observed future installations.}
\end{table}

The phrase ``3--7\% sensitivity'' means \emph{repeating the calculation with
different assumed annual growth rates}. It does not mean that the forecast has
a measured statistical uncertainty of 3--7\%, nor that a probability has been
assigned to the resulting range. ``Slow, central and fast growth scenarios'' is
the clearer presentation wording.

\section{Interpreting the central case}

With 5\% annual net growth, the first year's increase is approximately 15 rooms.
At the 2036 installed-base level, the following year's increase would be about
24 rooms; at the 2046 level, about 40 rooms. This describes a gradual increase
in annual deployment rather than a fixed number of new rooms every year.

For comparison, simply adding 19 rooms annually to the 300-room baseline would
give 490 rooms after ten years and 680 after twenty years, before allowing for
closures. The 19-room figure comes from the existing model's rounded estimate
of global 2025 room orders: IBA's 12 orders divided by its reported 63\% share.
An order is not a commissioned room; this comparison assumes a sustained
conversion of orders into deliveries~\cite{ibapresentation}.

The central case therefore implies a similar ten-year result to constant
additions at that rate, but some acceleration over twenty years. The fast case
requires considerably stronger expansion and should not be treated as the
default outcome.

The growth rate is defined as \emph{net} growth in operating rooms. It does not
separately predict closures, refurbishment or equipment replacement. Gross
new commissioning can exceed net growth when older rooms close. Replacement
of equipment in a continuing room does not increase the installed-room count.

\section{Sensitivity to the starting inventory}

The initial room count is a separate source of uncertainty from the growth
rate. Reducing the baseline from 300 to 250 lowers every projection by one sixth,
approximately 16.7\%.

\begin{table}[htbp]
\centering
\begin{tabular}{lrrr}
\toprule
Central case: 5\% annual growth & 2026 rooms & 2036 rooms & 2046 rooms\\
\midrule
Lower starting inventory & 250 & 407 & 663\\
Working starting inventory & 300 & 489 & 796\\
\bottomrule
\end{tabular}
\caption{Effect of changing the assumed starting inventory. Neither starting
value is presented here as an audited count.}
\end{table}

Verifying the baseline should therefore precede any attempt to refine the
forecast to narrow numerical ranges. A practical inventory should record the
facility, country, clinical proton rooms, operating status, opening date and
ARGOS compatibility. Construction and planning lists should be reconciled to
remove duplicates and projects that have already opened, closed or stalled.

\section{Implications for the ARGOS market}

The central case adds approximately 189 rooms by 2036 and 496 rooms by 2046
relative to the assumed 2026 baseline. At the current proposed ARGOS system
price of \EUR\,2.5 million, the corresponding net-expansion proxy represents
approximately \textbf{\EUR\,0.47 billion over ten years and \EUR\,1.24 billion
over twenty years}, before applying compatibility, monitoring adoption or
ARGOS market share. Values use the unrounded growth calculation and constant
2026 prices.

These amounts are cumulative potential system sales associated with net room
expansion, not annual revenues and not a forecast of ARGOS sales. They exclude
retrofits of the existing base and replacement demand. Because closures have
not been modelled separately, the net expansion is not an estimate of gross
new-build deliveries.

For an annual commercial forecast, the relevant relationship is
\begin{equation}
\begin{aligned}
\text{ARGOS system revenue}
&=\text{eligible installations}\times\text{monitoring adoption}\\
&\quad\times\text{ARGOS share}\times\text{system price}.
\end{aligned}
\end{equation}
The eligible-installation term must be constructed separately for new rooms
and retrofits. Service revenue is then determined by the active paying ARGOS
contract base, using the proposed \EUR\,250,000 annual fee, rather than by
multiplying every proton room by that fee.

\section{Recommended use in the investment discussion}

Use \textbf{5\% annual net room growth as the central working assumption}, with
3\% and 7\% cases to show the effect of slower and faster expansion. The rounded
central message is ``approximately 500 treatment rooms in ten years and 800 in
twenty years'', explicitly conditional on a 300-room starting inventory.

Before incorporating the forecast into the financial tables or presentation,
verify the starting room inventory and compare the assumed growth rates with a
consistent historical series. Do not substitute revenue growth rates from
equipment-market reports for growth in physical treatment-room numbers.

This note is separate from the existing four-table financial model. It does
not modify its installation volumes, pricing assumptions or presentation slides.

\begin{thebibliography}{9}
\small
\bibitem{yap2026}
J.~Yap, A.~Steinberg, H.~Norman, K.~P.~Nesteruk and S.~Sheehy,
``Towards ultra-fast treatments: large energy acceptance beam delivery systems
and opportunities for proton beam therapy'',
\emph{Frontiers in Oncology} \textbf{16}, 1791102 (3 June 2026).
\href{https://doi.org/10.3389/fonc.2026.1791102}{doi:10.3389/fonc.2026.1791102}.
Facility counts and single-room fraction are reported in the introduction and
Figure~1, using PTCOG data from early 2026.

\bibitem{ptcogconstruction}
Particle Therapy Co-Operative Group,
\emph{Facilities under construction}, update June 2026.
\url{https://ptcog.online/facilities-under-construction/}.
Accessed 13 September 2026.

\bibitem{iba2025}
IBA, \emph{Full-year 2025 results}, 27 March 2026.
\href{https://www.iba-worldwide.com/sites/default/files/2026-03/260327-iba-fy2025-results-en.pdf}{Official results release}.
Proton Therapy business discussion: equipment backlog and room deliveries.

\bibitem{ibapresentation}
IBA, \emph{Full-year results 2025 presentation}, 27 March 2026, p.~13.
\href{https://www.iba-worldwide.com/sites/default/files/2026-03/iba-fy25-results-presentation-en.pdf}{Official results presentation}.
Reported IBA room orders and market share for 2025.
\end{thebibliography}

\end{document}
